\documentclass[10pt]{article}
\usepackage{template/template}
\usepackage{amsmath}
\NewsSetKeys{date={2026年11月10日},vol={028},title={每日有趣新闻},author={编辑部},note={待补编短刊：已核查新闻选题与来源；未完成正式版式、图像和PDF验收。}}
\begin{document}
\begin{documents}
\newsitem{南极臭氧洞：面积第七大，却不是最深}{科学深读｜11月9日｜NASA Earth Observatory}{NASA地球观测站在11月9日发布了2026年南极臭氧洞的年度分析：按9月7日至10月13日旺季平均面积计，今年在48年监测记录中排名第七大，但臭氧浓度的最低值并没有同样达到历史极端。所谓臭氧洞并非大气中出现了真正的空洞，而是将臭氧总柱浓度低于220多布森单位的区域定义为臭氧洞。多布森单位衡量的是把一根大气柱中的臭氧全部压缩到标准条件下所得的等效厚度。面积、深度和持续时间是三个不同指标，不应混为一谈。NASA提供的9月18日卫星图展示的是年度最大面积当日的空间分布，不能被误称为11月9日的即时卫星照片。\\par
为何臭氧洞在全球限制氯氟烃数十年后仍可能出现很大的年份？南极平流层的极涡温度与环流决定了臭氧破坏反应的气象条件。低温有利于极地平流层云形成，使含氯储库物质转化为更容易在春季阳光下破坏臭氧的活性形态。2026年较强的南极极涡恰好提供了有利条件。另一方面，长期趋势仍受1987年《蒙特利尔议定书》控制消耗臭氧层物质的政策影响。某一年面积偏大不能单独推翻恢复趋势，正如某一年面积偏小也不能宣称问题已彻底解决。NASA预计南极臭氧层恢复至1980年水平仍要到本世纪70年代。此处的恢复时间属于模型预测，不是已经观测到的事实。\\par
NASA文章提供基于NOAA-20卫星OMPS观测的原始数据可视化，图像注明制作者Lauren Dauphin、数据来自NASA Ozone Watch。本稿目前仅保存图像出处，没有取得可归档的实际图像文件，正式发行前必须补齐并核查使用许可。\cite{nasaOzone}}
\newsitem{AI算力缺电，数据中心开始寻找旧涡轮与退役喷气发动机}{科技与能源｜11月9日｜路透社}{人工智能基础设施扩张使电力供应成为与GPU采购同样实际的工程约束。路透社11月9日调查称，全球燃气轮机订单超过100吉瓦，而年度生产能力估计只有60至70吉瓦，美国新设备交付周期平均约六年。部分开发商因此寻找二手或翻新涡轮，甚至将退役航空发动机改造成发电设备。\\par
这一市场现象的关键不在于旧设备必然更划算，而在于建设时间。大型数据中心若等待全新燃机交付，机房、服务器与电力基础设施的投资可能无法同步投产。翻新设备能够缩短部分等待时间，却带来剩余寿命、连续运行可靠性、维护历史和排放许可等不确定性。喷气发动机本来面向航空工况，不能直接假定其在地面长期稳定发电时拥有相同效率和维护周期。\\par
更深层的问题是，电力设备短缺只是系统瓶颈之一。电网接入、变电站建设、天然气供给和环保审批仍可能推迟项目。新增装机也不等于新增可用电量：发电功率是某时刻的供给能力，电量则还取决于运行时间。对投资者而言，AI推理需求增长与算力项目能够按时盈利之间仍有漫长的工程链条。本报道引用的产能数字为路透社行业调查口径，并非全球设备厂商统一发布的审计总量。\cite{reutersTurbines}}
\newsitem{中国延长部分稀土出口管制暂停期限}{经济政策｜11月9日｜路透社}{路透社11月9日报道，中国宣布将部分稀土出口管制措施的暂停期限延长两个月，至2027年1月10日。此前的管制与中美经贸磋商相关。应特别区分管制规则被废止、规则仍存在但暂缓实施、以及具体企业已取得出口许可这三种法律状态。暂停并不自动等于无限制出口，也不能仅凭此推断中美已经达成全面长期贸易协议。现有报道没有足够细节说明所有适用品类及审批环节的实际变化，后续应以商务部正式公告与海关执行口径为准。\cite{reutersRareEarth}}
\newsitem{COP31开幕：电气化目标与低碳电源不是同一件事}{气候政策｜11月9日｜路透社}{第31届联合国气候变化大会在土耳其安塔利亚开幕。会议主席国推动到2035年将全球经济超过三分之一的活动电气化的倡议。电气化可使终端设备脱离直接燃烧化石燃料，但若电力仍主要来自高碳电源，排放削减幅度便取决于发电结构与系统效率。该倡议尚不是各国已共同履行的强制性减排承诺。会议还面临能源价格、战争和AI数据中心用电增长带来的政策冲突。\cite{reutersCOP}}
\newsitem{数学拾趣：Lean中的数学家控制权}{数学与形式化｜11月6日预印本，11月10日检索}{Patel等九位作者在arXiv预印本《Lea: An Agent Backbone for Mathematician-Controlled Formalization in Lean》中提出面向数学家而非完全自主代理的形式化框架。它将运行、会话与类型化事件组织成API，并集成Lean工具、Loogle与验证机制。作者报告了在Durfee方块生成函数案例中，由数学家给出证明分解、系统协助完成Lean开发的经历。此项工作是工具与案例报告，并非发现新的数学定理。\\par
形式化证明的困难并不只在于找到策略：必须确保自然语言命题翻译正确、依赖库版本明确、最终代码通过可信内核检查。作者称系统能保留数学家的高层控制，但一个成功案例不能推导出普遍成功率。论文于11月6日提交arXiv；目前按预印本处理，未查得同行评议接受信息，也未在本刊独立运行其代码。因未充分核实11月9日当天所有数学新稿，本期不将此旧两日论文冒充当日新发表成果。\cite{lea}}
\newsitem{编辑说明}{发行状态}{本稿是待补编短刊，不满足完整深度稿数量、实际新闻图片归档、Actions构建和PDF视觉验收要求。不得作为正式第028期发行，也不得宣称邮件已附PDF送达。}
\bibliographystyle{plain}
\bibliography{refs}
\end{documents}
\end{document}